\documentclass[12pt,letterpaper]{article}
\usepackage[T5,T1]{fontenc}\usepackage{mlmodern}
\usepackage{amsmath}
\usepackage[polish,vietnamese,french,main=english]{babel}
\usepackage[expansion=false,protrusion]{microtype}
\title{The MLModern fonts}
\author{Daniel Benjamin Miller}
\date{Version 1.2\\January 18, 2021}
\usepackage{titlesec}\titleformat*{\section}{\bfseries}\titleformat*{\subsection}{\bfseries}
\usepackage{hologo}
\begin{document}
\maketitle
MLModern is a text and math font family with \hologo{(La)TeX} support, based on the design of Donald Knuth's Computer Modern and the Latin Modern project. Some find the default vector version of Computer Modern used by default in most \hologo{TeX} distributions to be spindly, sometimes making it hard to read on screen as well as on paper; this is in contrast with the older bitmap versions of Computer Modern. MLModern provides a sturdy rendition of the Computer Modern design.

MLModern is based on Latin Modern, a version of Computer Modern with enhanced multilingual support. For basic use, add \texttt{\textbackslash usepackage\{mlmodern\}} to your preamble, or (for non-\hologo{LaTeX} use) replace references to \texttt{lm*} fonts with those to \texttt{mlm*} ones. You may also replace the Computer Modern and Latin Modern fonts (as well as the Vn\TeX\ fonts, etc.) without modifying a \hologo{TeX} file (or changing metrics) by using the map file \texttt{mlm-substitute.map}.

You must have the Latin Modern fonts installed, since this package expects their encoding files to be present. MLModern is available in unhinted Type 1 format, and should be good for both printing and on-screen reading. OpenType support is planned for a future version.

\section*{Licensing information}

The MLModern fonts are
\begin{quote}
Copyright © 2003--2009 by B.\ Jackowski and J.\ M.\ Nowacki.\\
Copyright © 2021 by Daniel Benjamin Miller.
\end{quote}
and based on the original ``Computer Modern'' family by Donald Knuth. A script by Chuanren Wu was used in the creation of the original MLModern font set before later adjustments. MLModern is made available under the \hologo{LaTeX} Project Public License, version 1.3c or later; the full terms are in \texttt{LICENSE}. This work has the \textsc{lppl} maintenance status ``maintained.'' The current maintainer of this work is Daniel Benjamin Miller.

\section*{Sample}

\begin{otherlanguage}{french}
Frères humains, qui après nous vivez,\\
N'ayez les cueurs contre nous endurciz,\\
Car, si pitié de nous pouvres avez,\\
Dieu en aura plustost de vous merciz.\\
Vous nous voyez cy attachez cinq, six:

\noindent Quant de la chair, que trop avons nourrie,\\
Elle est pieça devorée et pourrie,\\
Et nous, les os, devenons cendre et pouldre.\\
De nostre mal personne ne s'en rie,\\
Mais priez Dieu que tous nous vueille absouldre!\\

\noindent \textsc{François Villon}, \frquote{Ballade des pendus}\\

\end{otherlanguage}

\begin{otherlanguage}{vietnamese}
Tiếng Việt là ngôn ngữ của người Việt và là ngôn ngữ chính thức tại Việt Nam. Đây là tiếng mẹ đẻ của khoảng 85\% dân cư Việt Nam cùng với hơn 4 triệu Việt kiều.\end{otherlanguage}\ (\textsc{Wikipedia})\footnote{Full optical size support, as in Knuth's Computer Modern. And now here's some Polish: \foreignlanguage{polish}{Język polski, polszczyzna – język z grupy zachodniosłowiańskiej (do której należą również czeski, kaszubski słowacki i języki łużyckie), stanowiącej część rodziny indoeuropejskiej. Funkcjonuje jako język urzędowy Polski oraz należy do oficjalnych języków Unii Europejskiej.} (\textsc{Wikipedia})}

\begin{align}
E_0 &= mc^2 \\
E &= \frac{mc^2}{\sqrt{1-\frac{v^2}{c^2}}}
\end{align}

\end{document}
